\begin{table}[htbp]
\centering
\caption{社群媒體使用程度與線上、線下自我認同之迴歸結果}
\label{tab:ext}
\footnotesize
\setlength{\tabcolsep}{8pt}
\renewcommand{\arraystretch}{1.12}
\begin{tabular}{l cccc}
\toprule
& \multicolumn{2}{c}{依變數：線上自我認同} & \multicolumn{2}{c}{依變數：線下自我認同} \\
\cmidrule(lr){2-3} \cmidrule(lr){4-5}
& 未調整 & 調整後 & 未調整 & 調整後 \\
\midrule
社群媒體使用程度 & $0.0268^{***}$ & $0.0275^{***}$ & $-0.0156^{***}$ & $-0.0154^{***}$ \\
 & (0.0015) & (0.0016) & (0.0015) & (0.0017) \\[1pt]
另一種自我認同分數 & $0.374^{***}$ & $0.374^{***}$ & $0.384^{***}$ & $0.382^{***}$ \\
 & (0.013) & (0.013) & (0.013) & (0.014) \\[1pt]
\midrule
背景控制變數 & 否 & 是 & 否 & 是 \\
預測變化（P5 至 P95） & -- & $+0.50$ & -- & $-0.28$ \\
$N$ & 8,781 & 7,948 & 8,781 & 7,948 \\
$R^2$ & .168 & .181 & .145 & .159 \\
\bottomrule
\end{tabular}

\vspace{0.4em}
\parbox{\linewidth}{\scriptsize 註：括號內為 HC3 穩健標準誤；社群媒體使用程度之 $p$ 值經 FDR 校正。「另一種自我認同分數」在線上認同模型中為線下認同，在線下認同模型中為線上認同。預測變化為使用程度由第 5 百分位提高至第 95 百分位時，依樣本實際分布平均之預測值變化。$^{*}p<.05$，$^{**}p<.01$，$^{***}p<.001$。}
\end{table}
